% Fragmento ES para inserción, no documento autónomo.
% Destino: V/main.tex, después de v_transporte_biparticion.tex y antes
% de 60_correlacion_cuasilibre_entropias.tex. ES, con pruebas contiguas.
\subsection{Energía recuperable dentro de una distribución de área fija}
\label{ins29:v:orientacion}

La frontera trítica permite separar el contenido de la ocupación total
y el de su procedencia sectorial. Recibimos los dos conjuntos de
dieciocho enlaces, sus ocupaciones \(0,1,2\), los pesos \(90/120\)
y la sección térmica construida anteriormente. Escribimos
\[
 X=N_{90},\qquad Y=N_{120},\qquad N=X+Y,\quad M=X-Y,
 \qquad d=30\Delta_{\rm mod}>0.
\]
Sea \(\epsilon=k_BT>0\) la escala energética de la realización
declarada. Su Hamiltoniano y su estado de referencia son
\[
 H=\epsilon d(3X+4Y),\qquad
 \rho_0=Z^{-1}e^{-H/\epsilon}.
\]
El operador de área ya construido es
\(\widehat{\mathcal A}=a_{\rm b}N\), con
\(a_{\rm b}=\gamma a_0L_\alpha^2>0\) fijo durante el protocolo.
Definimos \(W\) como el intercambio de los dos conjuntos de enlaces,
conservando cada ocupación. Esta permutación invierte \(M\) y deja
\(N\) invariante. La inversión complementaria de ocupaciones tratada
anteriormente, \(n\mapsto2-n\), conserva su función propia y es una
operación distinta.

\begin{proposition}[Respuesta energética invisible para la lectura areal]
\label{ins29:v:respuesta}
El estado \(\rho_1=W\rho_0W^*\) tiene la misma distribución completa
de área y la misma entropía de von Neumann que \(\rho_0\), mientras
\[
 \Delta E=\operatorname{Tr}[H(\rho_1-\rho_0)]
 =\epsilon\mathcal J(d)>0,
\]
\[
 \mathcal J(d)=D(\rho_1\Vert\rho_0)
 =18d\,[\mu(3d)-\mu(4d)],\qquad
 \mu(a)=\frac{e^{-a}+2e^{-2a}}{1+e^{-a}+e^{-2a}}.
\]
El intercambio inverso recupera exactamente \(\Delta E\); esta cantidad
es la máxima energía recuperable mediante una transformación unitaria
del estado \(\rho_1\), para el Hamiltoniano fijo \(H\).
\end{proposition}
\begin{proof}
Como \(WNW^*=N\), todos los proyectores espectrales del área conmutan
con \(W\) y conservan sus probabilidades. La unitariedad conserva el
espectro del estado y su entropía. En la base de ocupaciones, escribamos
\(p_0,p_1\) para las probabilidades. El intercambio da
\[
 \log\frac{p_1}{p_0}=-dM,\qquad
 \langle M\rangle_{p_1}=-\langle M\rangle_{p_0},\qquad
 \langle M\rangle_{p_0}=18[\mu(3d)-\mu(4d)].
\]
La suma ponderada del logaritmo prueba la fórmula de
\(D(\rho_1\Vert\rho_0)\). Además,
\(\mu'(a)=-\operatorname{Var}_a(n)<0\) para los tres pesos positivos,
por lo que \(\mathcal J(d)>0\).
La identidad de Gibbs, junto con la igualdad de entropías, da
\(\Delta E=\epsilon D(\rho_1\Vert\rho_0)\).
Para cualquier unitario \(V\), el estado \(V\rho_1V^*\) conserva esa
misma entropía y satisface
\[
 \operatorname{Tr}[H(V\rho_1V^*-\rho_0)]
 =\epsilon D(V\rho_1V^*\Vert\rho_0)\geq0.
\]
La elección \(V=W^*\) alcanza la igualdad y prueba el máximo.
\end{proof}

La distribución areal conserva todas las multiplicidades de \(N\),
pero el Hamiltoniano distingue sus repartos \((X,Y)\). El exceso
energético queda situado exactamente en esa distinción. La escala de
horizonte, cuando se utiliza la realización correspondiente, es
\(\epsilon=\hbar c/(2\pi R)\). Así, la acción convierte
\(\mathcal J(d)\) en una respuesta energética a radio fijo. El
operador de área de esta preparación y el área geométrica de un horizonte
con retroacción conservan sus dominios respectivos.

\subsubsection{Restitución condicionada y funcional reducido}

La localización intrafibra admite una formulación exacta. Sobre el
espacio finito de configuraciones, sea \(q\) un extractor sobreyectivo,
\(p_{0,q}=q_*p_0\) y \(p_q=q_*p\). Para una distribución \(r\) del
extractor, definimos
\[
 (\mathcal R_0r)(x)
 =r(q(x))\frac{p_0(x)}{p_{0,q}(q(x))}.
\]
Este restituidor satisface \(q_*\mathcal R_0=I\). Su distribución
condicionada conserva los pesos térmicos dentro de cada fibra.

\begin{proposition}[Energía libre de la información intrafibra]
\label{ins29:v:restitucion}
Para \(F_T(p)=\sum_xp(x)H(x)+\epsilon\sum_xp(x)\log p(x)\),
\[
 F_T(p)-F_T(\mathcal R_0p_q)
 =\epsilon D(p\Vert\mathcal R_0p_q).
\]
Por tanto, \(\mathcal R_0p_q\) es el único minimizador de la energía
libre entre las distribuciones con marginal \(p_q\).
\end{proposition}
\begin{proof}
En el soporte de \(p\), la factorización
\[
 \frac{p(x)}{p_0(x)}=
 \frac{p_q(q(x))}{p_{0,q}(q(x))}
 \frac{p(x)}{(\mathcal R_0p_q)(x)}
\]
da, al tomar logaritmos y sumar,
\(D(p\Vert p_0)=D(p_q\Vert p_{0,q})+
D(p\Vert\mathcal R_0p_q)\).
La identidad \(F_T(p)=\epsilon D(p\Vert p_0)-\epsilon\log Z\)
y la misma fórmula aplicada a \(\mathcal R_0p_q\) prueban el
enunciado. La divergencia se anula exactamente cuando sus argumentos
coinciden; se adopta \(0\log0=0\).
\end{proof}

Para \(q=N\), escribamos
\(g_{18}(x)=[t^x](1+t+t^2)^{18}\), la multiplicidad trítica ya
construida. El Hamiltoniano efectivo de la lectura areal es
\[
 \frac{H_N(n)}\epsilon=
 \frac{7d}{2}n-\log\!\sum_{x+y=n}
 g_{18}(x)g_{18}(y)e^{d(x-y)/2}.
\]
Esta expresión procede de sumar \(e^{-H/\epsilon}\) sobre la fibra.
El sumatorio es invariante bajo invertir el signo del término orientado,
por intercambio de \(x\) e \(y\). Para la preparación anterior,
\(p_{1,N}=p_{0,N}\) y \(\mathcal R_0p_{1,N}=p_0\): todo el exceso
de energía libre reside en la información que esa lectura areal omite.
En cambio, conservar \((N,M)\) fija la energía de cada fibra y hace
uniforme su distribución condicionada. El resultado precisa qué
observables bastan para recuperar esta colección y su respuesta térmica.
